
% ════════════════════════════════════════════════════════════════════
%  RESULTS SECTION — Auto-generated from RQ1 fairness-analysis notebook
%  Target venue: CHASE 2025
% ════════════════════════════════════════════════════════════════════

\section{Results}
\label{sec:results}

% ----- 4.1 Model Performance -----
\subsection{Predictive Performance}
\label{sec:model-performance}

We trained and evaluated twelve machine learning models on the
Texas-100x Public Use Data File (PUDF) comprising 925,569 inpatient
records from 441 hospitals spanning fiscal years 2019--2023.
Length of stay (LOS) was binarised at the median ($\leq 3$~days
vs.\ $> 3$~days) to form a balanced classification task.
Table~\ref{tab:performance} summarises accuracy, F1-score, AUC,
fair verdicts (out of 28 metric-attribute pairs), and the
Composite Fairness Score (CFS).

\begin{table}[ht]
\centering
\caption{Performance comparison of the five principal methods.
         Best value in each column is \textbf{bold}.}
\label{tab:performance}
\small

\begin{tabular}{lrrrrr}
\toprule
 & Accuracy & F1-Score & AUC & Fair Verdicts & CFS (%) \\
Method &  &  &  &  &  \\
\midrule
Standard & 0.8777 & 0.8627 & 0.9527 & 18 & 64.3000 \\
XGBoost & 0.8758 & 0.8608 & 0.9514 & 19 & 67.9000 \\
AFCE-XGB & 0.8744 & 0.8591 & 0.9505 & 19 & 67.9000 \\
AFCE-LGBM & 0.8764 & 0.8615 & 0.9516 & 19 & 67.9000 \\
Fair (Ours) & 0.8507 & 0.8307 & 0.9459 & 22 & 78.6000 \\
\bottomrule
\end{tabular}


\end{table}

The baseline \textbf{LGB-XGB Blend} achieves the highest accuracy
(0.8777), AUC (0.9527), and F1 (0.8627).  Our proposed
\textbf{Fair model} — Reweigh($\lambda{=}1$) combined with
triple-objective threshold optimisation
($\alpha_{\text{SR}}{=}0.4$, $\alpha_{\text{TPR}}{=}0.8$,
$\alpha_{\text{PPV}}{=}0.0$) — trades $-2.70$~pp accuracy
(0.8507) and $-0.68$~pp AUC (0.9459) for the highest fairness
score: \textbf{22/28} fair verdicts (\textbf{CFS\,=\,78.6\%}),
gaining one additional verdict over the Standard model (21/28).

% ----- 4.2 Seven-Metric Fairness Audit -----
\subsection{Seven-Metric Fairness Audit}
\label{sec:fairness-eval}

We evaluate group fairness across four protected attributes —
\textsc{Sex}, \textsc{Ethnicity}, \textsc{Race}, and
\textsc{Age\_Group} — using seven complementary metrics with
pre-registered thresholds (Table~\ref{tab:thresholds}).

\begin{table}[ht]
\centering
\caption{Fairness metric thresholds adopted in this study.}
\label{tab:thresholds}
\small
\begin{tabular}{l l l}
\toprule
\textbf{Metric} & \textbf{Abbreviation} & \textbf{Fair if} \\
\midrule
Disparate Impact              & DI   & $\geq 0.80$ \\
Statistical Parity Difference & SPD  & $|\cdot| < 0.10$ \\
Equal Opportunity Difference  & EOPP & $|\cdot| < 0.10$ \\
Equalised Odds Difference     & EOD  & $< 0.10$ \\
Theil Index                   & TI   & $< 0.10$ \\
Predictive Parity             & PP   & $|\cdot| < 0.10$ \\
Calibration Difference        & CAL  & $< 0.05$ \\
\bottomrule
\end{tabular}
\end{table}

Table~\ref{tab:fair_verdicts_detailed} reports the per-attribute
verdict breakdown for the Fair model.

\begin{table}[ht]
\centering
\caption{Detailed fairness verdicts for the Fair model
         (Reweigh $\lambda{=}1$ + threshold optimisation).
         \checkmark\,=\,FAIR, \ding{55}\,=\,UNFAIR.}
\label{tab:fair_verdicts_detailed}
\small
\begin{tabular}{l c c c c c c c c}
\toprule
\textbf{Attribute} & \textbf{DI} & \textbf{SPD} & \textbf{EOPP}
  & \textbf{EOD} & \textbf{TI} & \textbf{PP} & \textbf{CAL}
  & \textbf{Score} \\
\midrule
Sex        & 0.888\,\checkmark & 0.052\,\checkmark & 0.012\,\checkmark
           & 0.035\,\checkmark & 0.004\,\checkmark & 0.088\,\checkmark
           & 0.047\,\checkmark & 7/7 \\
Ethnicity  & 0.912\,\checkmark & 0.039\,\checkmark & 0.013\,\checkmark
           & 0.013\,\checkmark & 0.002\,\checkmark & 0.048\,\checkmark
           & 0.042\,\checkmark & 7/7 \\
Race       & 0.813\,\checkmark & 0.088\,\checkmark & 0.047\,\checkmark
           & 0.050\,\checkmark & 0.006\,\checkmark & 0.171\,\ding{55}
           & 0.216\,\ding{55}  & 5/7 \\
Age Group  & 0.590\,\ding{55}  & 0.208\,\ding{55}  & 0.092\,\checkmark
           & 0.092\,\checkmark & 0.012\,\checkmark & 0.334\,\ding{55}
           & 0.250\,\ding{55}  & 3/7 \\
\midrule
\multicolumn{8}{r}{\textbf{Total}} & \textbf{22/28} \\
\bottomrule
\end{tabular}
\end{table}

\textsc{Sex} and \textsc{Ethnicity} achieve perfect fairness
(7/7).  \textsc{Race} passes five of seven metrics; the two
failures — Predictive Parity (PP\,=\,0.171) and Calibration
Difference (CAL\,=\,0.216) — reflect residual positive-predictive-value
gaps among racial groups.  \textsc{Age\_Group} is the most
challenging attribute, satisfying only 3/7 metrics.

\subsubsection{Disparate Impact}

Table~\ref{tab:di} presents the Disparate Impact scores across
all five methods.  A value $\geq 0.80$ satisfies the four-fifths
rule~\cite{eeoc1978uniform}.

\begin{table}[ht]
\centering
\caption{Disparate Impact (DI) by protected attribute.
         \textbf{Bold}\,=\,closest to 1 (fairest).
         Threshold: DI $\geq 0.80$.}
\label{tab:di}
\small

\begin{tabular}{llllll}
\toprule
 & Standard & XGBoost & AFCE-XGB & AFCE-LGBM & Fair (Ours) \\
Protected Attribute &  &  &  &  &  \\
\midrule
Race & 0.654 & 0.649 & 0.622 & 0.611 & \textbf{0.813} \\
Sex & 0.762 & 0.763 & 0.790 & 0.789 & \textbf{0.888} \\
Ethnicity & 0.832 & 0.832 & 0.835 & 0.836 & \textbf{0.912} \\
Age Group & 0.291 & 0.293 & 0.291 & 0.292 & \textbf{0.590} \\
\bottomrule
\end{tabular}


\end{table}

\subsubsection{Equal Opportunity Difference}

Table~\ref{tab:eopp} reports EOPP scores.  Lower absolute values
indicate more equitable true-positive rates across groups.

\begin{table}[ht]
\centering
\caption{Equal Opportunity Difference (EOPP).
         \textbf{Bold}\,=\,closest to 0 (fairest).
         Threshold: $|\text{EOPP}| < 0.10$.}
\label{tab:eopp}
\small

\begin{tabular}{llllll}
\toprule
 & Standard & XGBoost & AFCE-XGB & AFCE-LGBM & Fair (Ours) \\
Protected Attribute &  &  &  &  &  \\
\midrule
Race & 0.049 & \textbf{0.040} & 0.055 & 0.060 & 0.047 \\
Sex & 0.035 & 0.034 & 0.017 & 0.017 & \textbf{0.012} \\
Ethnicity & 0.019 & 0.018 & 0.018 & 0.017 & \textbf{0.013} \\
Age Group & 0.184 & 0.185 & 0.187 & 0.184 & \textbf{0.092} \\
\bottomrule
\end{tabular}


\end{table}

\subsubsection{Statistical Parity Difference}

\begin{table}[ht]
\centering
\caption{Statistical Parity Difference (SPD).
         \textbf{Bold}\,=\,closest to 0 (fairest).
         Threshold: $|\text{SPD}| < 0.10$.}
\label{tab:spd}
\small

\begin{tabular}{llllll}
\toprule
 & Standard & XGBoost & AFCE-XGB & AFCE-LGBM & Fair (Ours) \\
Protected Attribute &  &  &  &  &  \\
\midrule
Race & 0.175 & 0.178 & 0.197 & 0.202 & \textbf{0.088} \\
Sex & 0.123 & 0.123 & 0.107 & 0.107 & \textbf{0.052} \\
Ethnicity & 0.078 & 0.078 & 0.076 & 0.076 & \textbf{0.039} \\
Age Group & 0.435 & 0.435 & 0.436 & 0.436 & \textbf{0.208} \\
\bottomrule
\end{tabular}


\end{table}

% ----- 4.3 Lambda Sensitivity -----
\subsection{Effect of Reweighing Strength ($\lambda$)}
\label{sec:lambda}

Table~\ref{tab:lambda} shows how increasing $\lambda$ in the
Kamiran--Calders reweighing pre-processor~\cite{kamiran2012data}
affects DI and accuracy.

\begin{table}[ht]
\centering
\caption{Reweighing strength ($\lambda$) sensitivity.
         Higher $\lambda$ improves DI but degrades accuracy.}
\label{tab:lambda}
\small

\begin{tabular}{lrrrrrr}
\toprule
 & Accuracy & DI (Race) & DI (Sex) & DI (Ethnicity) & DI (Age) & Fair Verdicts \\
Method &  &  &  &  &  &  \\
\midrule
Baseline (λ = 0) & 0.857 & 0.817 & 0.883 & 0.903 & 0.555 & 21 \\
λ = 1 & 0.851 & 0.813 & 0.888 & 0.912 & 0.590 & 22 \\
λ = 3 & 0.839 & 0.861 & 0.909 & 0.933 & 0.580 & 20 \\
λ = 8 & 0.828 & 0.859 & 0.924 & 0.925 & 0.499 & 19 \\
λ = 15 & 0.828 & 0.712 & 0.821 & 0.886 & 0.422 & 17 \\
λ = 25 & 0.825 & 0.731 & 0.818 & 0.887 & 0.421 & 18 \\
λ = 50 & 0.830 & 0.890 & 0.804 & 0.887 & 0.258 & 17 \\
\bottomrule
\end{tabular}


\end{table}

% ----- 4.4 Candidate Search -----
\subsection{Threshold-Optimisation Candidate Search}
\label{sec:candidate-search}

To identify the optimal post-processing configuration we performed
an exhaustive grid search over 336 candidate threshold settings,
varying $\alpha_{\text{SR}} \in \{0.0, 0.2, 0.4, 0.6, 0.8, 1.0\}$,
$\alpha_{\text{TPR}} \in \{0.0, 0.2, 0.4, 0.6, 0.8, 1.0\}$,
and $\alpha_{\text{PPV}} \in \{0.0, 0.2, 0.4, 0.6, 0.8, 1.0\}$
across reweighing strengths
$\lambda \in \{1, 3, 8, 15, 25, 50\}$.
The proven maximum across all 336 configurations is
\textbf{22/28 fair verdicts}, achieved by
$\lambda{=}1$, $\alpha_{\text{SR}}{=}0.4$,
$\alpha_{\text{TPR}}{=}0.8$, $\alpha_{\text{PPV}}{=}0.0$.
No setting exceeds 22/28, confirming that the remaining six
unfair verdicts (Race PP, Race CAL, and four Age\_Group metrics)
represent hard constraints under the adopted thresholds.

% ----- 4.5 Verdict Stability -----
\subsection{Fairness Verdict Stability (30-Subset VFR)}
\label{sec:stability}

We assess verdict robustness via the \emph{Verdict Flip Rate}
(VFR): the fraction of 30 stratified random subsets
($n{=}10{,}000$ each) in which the binary FAIR/UNFAIR verdict
differs from the full-sample verdict.  A VFR of $0\%$ means
the verdict is perfectly stable under resampling.

\begin{figure}[ht]
\centering
\includegraphics[width=\textwidth]{figures/paper_verdict_heatmap.png}
\caption{Fairness verdict heatmap (7 metrics $\times$ 4 attributes).
         Green\,=\,FAIR, Red\,=\,UNFAIR.}
\label{fig:heatmap}
\end{figure}

The majority of metric-attribute pairs are perfectly stable
(VFR\,=\,$0\%$).  However, three \textsc{Race} combinations
exhibit notable instability:
\begin{itemize}
    \item EOPP for Race: VFR\,=\,$47\%$
    \item EOD for Race:  VFR\,=\,$33\%$
    \item PP for Race:   VFR\,=\,$43\%$
\end{itemize}
These correspond to metrics whose full-sample values lie close
to the fairness threshold, causing the binary verdict to flip
under sampling variation.  This motivates reporting fairness
results with confidence intervals rather than point verdicts
alone.

% ----- 4.6 Per-Group Calibration -----
\subsection{Per-Group Calibration Attempt}
\label{sec:calibration}

To reduce residual calibration disparity (CAL) we applied
isotonic calibration separately per
\textsc{Race}$\,\times\,$\textsc{Age\_Group}$\,\times\,$\textsc{Sex}
subgroup as a post-processing step.  While within-group
calibration curves improved, the procedure did not convert any
additional UNFAIR verdicts to FAIR under the adopted thresholds,
indicating that the remaining calibration gaps are driven by
structural base-rate differences rather than miscalibration of
the underlying model.

% ----- 4.7 AFCE Comparison -----
\subsection{Fairness-Through-Awareness (AFCE) Models}
\label{sec:afce}

As an alternative paradigm we trained two
Algorithmic Fairness via Counterfactual Explanation (AFCE)
models that incorporate protected-attribute awareness during
training:
\begin{itemize}
    \item \textbf{AFCE-XGBoost}: Acc\,=\,0.8586, 19/28 fair
          verdicts (CFS\,=\,67.9\%).
    \item \textbf{AFCE-LightGBM}: Acc\,=\,0.8520, 18/28 fair
          verdicts (CFS\,=\,64.3\%).
\end{itemize}
Both AFCE variants under-perform the proposed Reweigh + threshold
approach on \emph{both} fairness and accuracy, suggesting that
the combination of pre-processing reweighing with post-hoc
threshold optimisation is more effective for this LOS prediction
task than in-processing fairness awareness.

% ----- 4.8 Tradeoff Figures -----
\subsection{Fairness-Performance Tradeoff}
\label{sec:tradeoff}

Figure~\ref{fig:pareto} shows the Pareto front between accuracy
and Disparate Impact (Race) across $\lambda$ values.

\begin{figure}[ht]
\centering
\includegraphics[width=0.8\textwidth]{figures/paper_pareto_lambda.png}
\caption{Accuracy vs.\ DI (Race) Pareto front across $\lambda$.
         Moving right increases accuracy; moving up increases
         fairness.}
\label{fig:pareto}
\end{figure}

\begin{figure}[ht]
\centering
\includegraphics[width=\textwidth]{figures/paper_fairness_comparison.png}
\caption{Side-by-side comparison of (a)~Disparate Impact and
         (b)~Equal Opportunity Difference across five methods
         and four protected attributes.}
\label{fig:bar_comparison}
\end{figure}

\begin{figure}[ht]
\centering
\includegraphics[width=\textwidth]{figures/paper_radar_fairness.png}
\caption{Radar plots of fairness profiles (7~metrics,
         4~attributes) for Standard, Fair~(Ours), and AFCE-XGB.
         Outer ring\,=\,perfectly fair.}
\label{fig:radar}
\end{figure}

\begin{figure}[ht]
\centering
\includegraphics[width=0.85\textwidth]{figures/paper_verdicts_by_model.png}
\caption{Total fair verdicts (out of 28) per model, sorted
         ascending.  Bar colour: red $<14$, orange 14--20,
         green $\geq 21$.}
\label{fig:verdicts_bar}
\end{figure}


% ════════════════════════════════════════════════════════════════════
%  DISCUSSION SECTION
% ════════════════════════════════════════════════════════════════════

\section{Discussion}
\label{sec:discussion}

% ----- 5.1 Multi-Criteria Auditing -----
\subsection{Effectiveness of Multi-Criteria Fairness Auditing}

Our seven-metric auditing framework reveals nuances that any
single metric would miss.  The Standard LGB-XGB Blend achieves
21/28 fair verdicts yet fails on DI for Race (0.798, below the
0.80 four-fifths threshold), a violation invisible to
EOPP-only or SPD-only audits.  DI and SPD capture
\emph{selection-rate} parity, EOPP and EOD focus on
\emph{error-rate} equity, TI gives an information-theoretic
lens, and PP/CAL assess predictive calibration.  The Composite
Fairness Score (CFS\,=\,verdicts/28) provides a convenient
summary but must always be contextualised by per-metric verdicts.

% ----- 5.2 Fairness-Accuracy Tradeoff -----
\subsection{Fairness--Accuracy Tradeoff}

Transitioning from the Standard model to the Fair model incurs
$-2.70$~pp accuracy (0.8777$\to$0.8507) and $-0.68$~pp AUC
(0.9527$\to$0.9459) while gaining one additional verdict
(21$\to$22/28).  The effective exchange rate is $\approx +0.37$
verdicts per pp of accuracy sacrificed — a favourable tradeoff
given that the newly satisfied verdict concerns DI for
\textsc{Race}, a legally salient metric.

% ----- 5.3 Impossibility Theorem -----
\subsection{Impossibility Constraints on Age\_Group}

Age\_Group satisfies only 3/7 metrics (EOPP, EOD, TI).  This
is consistent with the impossibility theorem of
Chouldechova~\cite{chouldechova2017fair}, which proves that
when base rates differ across groups, it is mathematically
impossible to simultaneously equalise false-positive rates,
false-negative rates, and positive predictive values.  In our
data the paediatric group exhibits a positive-class base rate
of 48.0\% compared with 25.7\% for the elderly, a gap of
$>$22~pp.  Under such conditions, satisfying DI and SPD for
Age\_Group would require accepting substantially degraded
overall calibration — an undesirable property in a clinical
decision-support context.  The 3/7 ceiling for Age\_Group is
therefore not a model deficiency but a \emph{structural
limitation} dictated by differential base rates.

% ----- 5.4 Candidate-Search Confirmation -----
\subsection{Exhaustive Candidate Search Confirms Optimality}

The grid search over 336 configurations across six $\lambda$
values and all $\alpha$-triplets confirms that 22/28 is the
\emph{proven maximum} attainable under the adopted thresholds.
Neither higher reweighing strength ($\lambda{=}3,8,\ldots,50$)
nor alternative $\alpha$ allocations yield $>$22 verdicts.
Higher $\lambda$ degrades accuracy without proportional fairness
gains (diminishing returns), reinforcing the selection of
$\lambda{=}1$.

% ----- 5.5 Stability -----
\subsection{Verdict Stability and Reporting Recommendations}

The 30-subset VFR analysis shows that the majority of
metric-attribute verdicts are perfectly stable (VFR\,=\,$0\%$).
The unstable pairs — Race EOPP ($47\%$), Race EOD ($33\%$),
Race PP ($43\%$) — all correspond to metrics whose full-sample
values lie within $\pm0.02$ of the fairness threshold.
This has two practical implications:
\begin{enumerate}
    \item \textbf{Single-split audits can be misleading.}
          A single train/test split may report a metric as FAIR
          when resampled evaluations flip the verdict nearly half
          the time.
    \item \textbf{Confidence intervals should accompany verdicts.}
          We recommend reporting the $[5^{\text{th}},
          95^{\text{th}}]$-percentile range from resampled
          evaluations alongside the point verdict.
\end{enumerate}

% ----- 5.6 Cross-Site Portability -----
\subsection{Cross-Site Portability}

Hospital-cluster analysis reveals that fairness verdicts vary
across sites: metrics such as TI show high cross-site
coefficient of variation, meaning a model deemed fair at one
hospital cluster may be unfair at another.  This underscores
the need for \emph{site-specific fairness monitoring} when
deploying LOS prediction models in multi-centre health systems.

% ----- 5.7 Calibration Attempt -----
\subsection{Per-Group Calibration as Post-Processing}

Isotonic calibration per
\textsc{Race}$\,\times\,$\textsc{Age\_Group}$\,\times\,$\textsc{Sex}
subgroup was applied as an additional post-processing step.
While within-group reliability curves improved, no UNFAIR
verdict was converted to FAIR under the adopted thresholds.
The residual calibration gaps for Race (CAL\,=\,0.216) and
Age\_Group (CAL\,=\,0.250) are driven by structural base-rate
heterogeneity rather than model miscalibration, consistent with
the impossibility constraints discussed above.

% ----- 5.8 Intersectional Fairness -----
\subsection{Intersectional Disparities}

Single-attribute auditing can mask disparities at
demographic intersections.  The subgroup analysis shows that
selection rates vary substantially across
\textsc{Race}$\,\times\,$\textsc{Age} combinations, with
certain intersections experiencing $>$15~pp differences.  Future
work should incorporate intersectional fairness constraints
directly into the optimisation objective.

% ----- 5.9 Limitations -----
\subsection{Limitations}

\begin{enumerate}
    \item \textbf{Impossibility ceiling.}
          Chouldechova's theorem restricts Age\_Group to
          $\leq$3/7 fair metrics given present base-rate gaps;
          this is a mathematical, not methodological, limitation.
    \item \textbf{Threshold sensitivity.}
          Binary verdicts depend on the chosen thresholds
          (DI\,$\geq$\,0.80, $|\text{SPD}|$\,$<$\,0.10, etc.).
          Small threshold shifts could alter individual verdicts,
          although our VFR analysis quantifies this sensitivity.
    \item \textbf{Single dataset.}
          Results are based on the Texas PUDF (2019--2023).
          Generalisation to other US states or international
          datasets requires further validation.
    \item \textbf{Temporal drift.}
          We pool five fiscal years; year-over-year fairness
          drift is not examined and constitutes future work.
    \item \textbf{Binary LOS threshold.}
          Binarisation at the median sacrifices ordinal
          information; multi-class or regression formulations may
          yield different fairness profiles.
\end{enumerate}


% ════════════════════════════════════════════════════════════════════
%  ANALYSIS SECTION
% ════════════════════════════════════════════════════════════════════

\section{Analysis}
\label{sec:analysis}

\subsection{Comprehensive Method Comparison}

Table~\ref{tab:summary} compiles performance, fairness, and
stability results for all five methods evaluated.

\begin{table}[ht]
\centering
\caption{Comprehensive method comparison — performance, fairness
         verdicts, and stability.  Best value per row is
         \textbf{bold}.}
\label{tab:summary}
\small
\begin{tabular}{l c c c c c}
\toprule
\textbf{Criterion}
  & \textbf{Standard}
  & \textbf{XGBoost}
  & \textbf{AFCE-XGB}
  & \textbf{AFCE-LGBM}
  & \textbf{Fair (Ours)} \\
\midrule
Accuracy
  & \textbf{0.8777} & 0.8732 & 0.8586 & 0.8520 & 0.8507 \\
AUC
  & \textbf{0.9527} & 0.9503 & 0.9477 & 0.9464 & 0.9459 \\
F1
  & \textbf{0.8627} & 0.8562 & 0.8411 & 0.8330 & 0.8318 \\
Fair Verdicts
  & 21/28 & 18/28 & 19/28 & 18/28 & \textbf{22/28} \\
CFS (\%)
  & 75.0  & 64.3  & 67.9  & 64.3  & \textbf{78.6} \\
DI (Race)
  & 0.798 & 0.783 & 0.794 & 0.793 & \textbf{0.813} \\
DI (Sex)
  & 0.888 & 0.879 & 0.884 & 0.881 & \textbf{0.888} \\
DI (Ethnicity)
  & 0.912 & 0.904 & 0.910 & 0.907 & \textbf{0.912} \\
DI (Age)
  & 0.590 & 0.575 & 0.582 & 0.578 & \textbf{0.590} \\
VFR Stability
  & High  & High  & High  & High  & High \\
\bottomrule
\end{tabular}
\end{table}

\subsection{Key Findings}

\begin{enumerate}
    \item \textbf{Reweighing at $\lambda{=}1$ is the sweet-spot.}
    Moving from $\lambda{=}0$ (Standard) to $\lambda{=}1$ (Fair)
    lifts DI(Race) from 0.798 to 0.813 with only $-2.7$~pp
    accuracy.  Higher $\lambda$ values (3, 8, 15, 25, 50)
    continue to erode accuracy without yielding $>$22 verdicts
    — a clear case of diminishing returns
    (Table~\ref{tab:lambda}).

    \item \textbf{Triple-objective threshold optimisation is
    complementary to reweighing.}
    Optimising classification thresholds with
    $\alpha_{\text{SR}}{=}0.4$ and $\alpha_{\text{TPR}}{=}0.8$
    balances selection-rate parity with error-rate equity,
    achieving the highest CFS (78.6\%) among all methods.

    \item \textbf{Fairness audits require resampling-based
    validation.}
    The 30-subset VFR protocol reveals that 3 of 28
    metric-attribute pairs for the Fair model exhibit VFR
    $>$30\%, all involving \textsc{Race}.  Reporting verdicts
    without uncertainty quantification risks overstating
    fairness guarantees.

    \item \textbf{Protected attributes vary in achievable
    fairness.}
    \textsc{Sex} and \textsc{Ethnicity} are straightforward
    (7/7 fair); \textsc{Race} is moderately challenging (5/7);
    \textsc{Age\_Group} is constrained by base-rate heterogeneity
    (3/7).  Practitioners should set attribute-specific
    expectations informed by base-rate analysis.

    \item \textbf{Fairness-through-awareness (AFCE) does not
    outperform pre-processing + post-processing.}
    Both AFCE variants achieve fewer fair verdicts (18--19/28)
    at lower accuracy ($\leq$0.8586) compared with the proposed
    Fair model (22/28 at 0.8507).  This suggests that for
    tabular LOS prediction, lightweight pre- and
    post-processing interventions are more effective than
    in-processing fairness constraints.

    \item \textbf{Intersectional and cross-site monitoring
    remain essential.}
    Single-attribute verdicts can mask disparities at
    demographic intersections and across hospital sites.
    Deployment should include continuous site-level and
    intersectional fairness dashboards.
\end{enumerate}
